\documentclass{beamer}
\usepackage[utf8]{inputenc}

\usetheme{Madrid}
\usecolortheme{default}
\usepackage{amsmath,amssymb,amsfonts,amsthm}
\usepackage{txfonts}
\usepackage{tkz-euclide}
\usepackage{listings}
\usepackage{adjustbox}
\usepackage{array}
\usepackage{tabularx}
\usepackage{gvv}
\usepackage{lmodern}
\usepackage{circuitikz}
\usepackage{tikz}
\usepackage[utf8]{inputenc}

\usetheme{Madrid}
\usecolortheme{default}
\usepackage{amsmath,amssymb,amsfonts,amsthm}
\usepackage{txfonts}
\usepackage{tkz-euclide}
\usepackage{listings}
\usepackage{adjustbox}
\usepackage{array}
\usepackage{tabularx}
\usepackage{gvv}
\usepackage{lmodern}
\usepackage{circuitikz}
\usepackage{tikz}
\usepackage{graphicx}

\setbeamertemplate{page number in head/foot}[totalframenumber]

\usepackage{tcolorbox}
\tcbuselibrary{minted,breakable,xparse,skins}



\definecolor{bg}{gray}{0.95}
\DeclareTCBListing{mintedbox}{O{}m!O{}}{%
  breakable=true,
  listing engine=minted,
  listing only,
  minted language=#2,
  minted style=default,
  minted options={%
    linenos,
    gobble=0,
    breaklines=true,
    breakafter=,,
    fontsize=\small,
    numbersep=8pt,
    #1},
  boxsep=0pt,
  left skip=0pt,
  right skip=0pt,
  left=25pt,
  right=0pt,
  top=3pt,
  bottom=3pt,
  arc=5pt,
  leftrule=0pt,
  rightrule=0pt,
  bottomrule=2pt,

  colback=bg,
  colframe=orange!70,
  enhanced,
  overlay={%
    \begin{tcbclipinterior}
    \fill[orange!20!white] (frame.south west) rectangle ([xshift=20pt]frame.north west);
    \end{tcbclipinterior}},
  #3,
}
\lstset{
    language=C,
    basicstyle=\ttfamily\small,
    keywordstyle=\color{blue},
    stringstyle=\color{orange},
    commentstyle=\color{green!60!black},
    numbers=left,
    numberstyle=\tiny\color{gray},
    breaklines=true,
    showstringspaces=false,
}
%------------------------------------------------------------
%This block of code defines the information to appear in the
%Title page
\title %optional
{2.3.17}
\date{August  2025}
%\subtitle{A short story}

\author % (optional)
{Bhoomika L - EE25BTECH11014}

\begin{document}


\frame{\titlepage}
\begin{frame}{Question}Find the properties of the triangle with vertices \brak{−4, 0},\brak{0, 3} and \brak{4, 0}.

\end{frame}

\begin{frame}{Theoretical Solution}
Given the points,
\begin{table}[H]    
  \centering
  \begin{tabular}[12pt]{ |c| c|}
    \hline
    \textbf{Name} & \textbf{Point}\\ 
    \hline
	$\vec{A}$ & \myvec{-4\\0}\\
    \hline 
	$\vec{B}$ & \myvec{0\\3}\\
    \hline
	$\vec{C}$ & \myvec{4\\0}\\
    \hline
    \end{tabular}
  \caption{Variables used.}
  \label{tab:2.3.17}
\end{table}
\end{frame}

\begin{frame}{Theoretical Solution}
We can calculate the magnitude of sides of the triangle using norms:
\begin{align}
    \brak{\vec{A}-\vec{B}}=\myvec{-4\\0}-\myvec{0\\3}=\myvec{-4\\3}
\end{align}
\begin{align}
 \brak{\vec{B}-\vec{C}}=\myvec{0\\3}-\myvec{4\\0}=\myvec{-4\\3}
 \end{align}
 \begin{align}
 \brak{\vec{A}-\vec{C}}=\myvec{-4\\0}-\myvec{4\\0}=\myvec{-8\\0}
 \end{align}
\end{frame}

\begin{frame}{Checking for orthogonality}
\begin{align}
 \brak{\vec{A}-\vec{B}}^T \cdot \brak{ \vec{B}-\vec{C}}=\myvec{-4&-3}\cdot \myvec{-4\\3}=16-9=7\neq 0\\
     \brak{ \vec{B}-\vec{C}}^T \cdot  \brak{\vec{A}-\vec{C}}= \myvec{-4&3} \cdot \myvec{-8\\0}=32+0=32 \neq 0\\
      \brak{\vec{A}-\vec{C}}^T \cdot \brak{\vec{A}-\vec{B}}=\myvec{-8&0} \cdot \myvec{-4\\-3}=32-0=32 \neq 0\\
\end{align}
 Hence the triangle is \textbf{not RIGHT} angled.
\end{frame}

\begin{frame}{To Check if the triangle is ISOSCELES}
If the triangle is ISOSCELES,a median should be perpendicular to a side of the triangle.\\
	Let the midpoints of the sides \textbf{AB}, \textbf{BC}, \textbf{AC} be $\vec{x}, \vec{y}$ and $\vec{z}$.
\begin{align}
\vec{x}=\frac{\vec{A}+\vec{B}}{2}=\myvec{-4\\0}+\myvec{0\\3}=\myvec{-2\\\frac{3}{2}}\\ \vec{y}=\frac{\vec{B}+\vec{C}}{2}=\myvec{0\\3}+\myvec{4\\0}=\myvec{2\\\frac{3}{2}}\\
\vec{z}=\frac{\vec{A}+\vec{C}}{2}=\myvec{-4\\0}+\myvec{4\\0}=\myvec{0\\0} 
\end{align}
\end{frame}

\begin{frame}{Theoretical Solution}
The vectors of medians of the triangle are:
\begin{align}
    \vec{x}-\vec{C}=\myvec{-2\\\frac{3}{2}}-\myvec{4\\0}=\myvec{-6\\\frac{3}{2}}\\
    \vec{y}-\vec{A}=\myvec{2\\\frac{3}{2}}-\myvec{-4\\0}=\myvec{6\\\frac{3}{2}}\\
    \vec{z}-\vec{B}=\myvec{0\\0}-\myvec{0\\3}=\myvec{0\\\frac{3}{2}}
\end{align}
\end{frame}

\begin{frame}{Theoretical Solution}
Then,
\begin{align}
  \brak{\vec{x}-\vec{C}}^T \cdot \brak{\vec{A}-\vec{B}}=\myvec{-6&\frac{3}{2}} \cdot \myvec{-4\\-3}=24-\frac{9}{2}=\frac{39}{2} \neq 0\\
     \brak{ \vec{y}-\vec{A}}^T \cdot \brak{\vec{B}-\vec{C}}= \myvec{6&\frac{3}{2}} \cdot \myvec{-4\\3}=-24+\frac{9}{2}=\frac{-39}{2} \neq 0\\
     \brak{ \vec{z}-\vec{B}}^T \cdot \brak{\vec{A}-\vec{C}}= \myvec{0&\frac{3}{2}} \cdot \myvec{-8\\0}=0
\end{align}
$ \therefore$ We can conclude that triangle \textbf{ABC} is an \textbf{ISOSCELES} triangle.
\end{frame}

\begin{frame}[fragile]
    \frametitle{C Code }
\begin{lstlisting}
#include <stdio.h>

// Function to check if triangle is right-angled
// Returns 1 if right-angled, 0 otherwise
int is_right_triangle() {
    // Vertices
    float Ax = -4, Ay = 0;
    float Bx = 0,  By = 3;
    float Cx = 4,  Cy = 0;

    // Side vectors
    float AB[2] = {Ax - Bx, Ay - By};
    float BC[2] = {Bx - Cx, By - Cy};

\end{lstlisting}
\end{frame}

\begin{frame}[fragile]
    \frametitle{C Code }
\begin{lstlisting}
   float AC[2] = {Ax - Cx, Ay - Cy};

    // Dot products
    float dot1 = AB[0]*BC[0] + AB[1]*BC[1];
    float dot2 = BC[0]*AC[0] + BC[1]*AC[1];
    float dot3 = AC[0]*AB[0] + AC[1]*AB[1];

    if (dot1 == 0 || dot2 == 0 || dot3 == 0)
        return 1;   // Right-angled
    else
        return 0;   // Not right-angled
}
// Function to check if triangle is isosceles
     // Returns 1 if isosceles, 0 otherwise
int is_isosceles() {
   

\end{lstlisting}
\end{frame}

\begin{frame}[fragile]
    \frametitle{C Code }
\begin{lstlisting}
    // Vertices
    float Ax = -4, Ay = 0;
    float Bx = 0,  By = 3;
    float Cx = 4,  Cy = 0;
    
    // Midpoints of sides
    float x_mid[2] = {(Ax + Bx) / 2, (Ay + By) / 2};
    float y_mid[2] = {(Bx + Cx) / 2, (By + Cy) / 2};
    float z_mid[2] = {(Ax + Cx) / 2, (Ay + Cy) / 2};

    
   // Medians
    float median1[2] = {x_mid[0] - Cx, x_mid[1] - Cy};
    float median2[2] = {y_mid[0] - Ax, y_mid[1] - Ay};
    float median3[2] = {z_mid[0] - Bx, z_mid[1] - By};
\end{lstlisting}
\end{frame}

\begin{frame}[fragile]
    \frametitle{C Code }
\begin{lstlisting}

 // Side vectors
    float AB[2] = {Ax - Bx, Ay - By};
    float BC[2] = {Bx - Cx, By - Cy};
    float AC[2] = {Ax - Cx, Ay - Cy};

 // Dot products (median $\perp$ opposite side ?)
    float dot1 = median1[0]*AB[0] + median1[1]*AB[1];
    float dot2 = median2[0]*BC[0] + median2[1]*BC[1];
    float dot3 = median3[0]*AC[0] + median3[1]*AC[1];

    if (dot1 == 0 || dot2 == 0 || dot3 == 0)
        return 1;  // Isosceles
    else
        return 0;  // Not isosceles}
\end{lstlisting}
\end{frame}

\begin{frame}[fragile]
    \frametitle{Python Code }
\begin{lstlisting}
import ctypes
import numpy as np
import matplotlib.pyplot as plt

# Load C library
c_lib = ctypes.CDLL('./triangle.so')

# Function signature
triangle_properties = c_lib.triangle_properties
triangle_properties.argtypes = [
    ctypes.POINTER(ctypes.c_double),  # sides[3]
    ctypes.POINTER(ctypes.c_double)   # angles[3]
\end{lstlisting}
\end{frame}


\begin{frame}[fragile]
    \frametitle{Python Code }
\begin{lstlisting}
]
triangle_properties.restype = None

# Prepare arrays
sides = (ctypes.c_double * 3)()
angles = (ctypes.c_double * 3)()

# Call the C function
triangle_properties(sides, angles)

# Convert to Python lists
sides_py = [sides[i] for i in range(3)]
angles_py = [angles[i] for i in range(3)]
\end{lstlisting}
\end{frame}

\begin{frame}[fragile]
    \frametitle{Python Code }
\begin{lstlisting}
 
print(f"Sides (AB, BC, CA): {sides_py}")
print(f"Angles (A, B, C): {angles_py} (degrees)")

# --- Plotting ---
A = np.array([-4, 0])
B = np.array([0, 3])
C = np.array([4, 0])

plt.figure(figsize=(8,6))
x_coords = [A[0], B[0], C[0], A[0]]
y_coords = [A[1], B[1], C[1], A[1]]

plt.plot(x_coords, y_coords, 'b-', linewidth=2, label="Triangle ABC")
\end{lstlisting}
\end{frame}


\begin{frame}[fragile]
\frametitle{Python Code}
\begin{lstlisting}
 # Mark vertices
plt.scatter([A[0], B[0], C[0]], [A[1], B[1], C[1]], color="red")
plt.text(A[0]-1, A[1]-0.3, f"A(-4,0)\n∠A≈{angles_py[0]:.2f}°", color="red")
plt.text(B[0]+0.2, B[1]+0.3, f"B(0,3)\n∠B≈{angles_py[1]:.2f}°", color="green")
plt.text(C[0]+0.4, C[1]-0.3, f"C(4,0)\n∠C≈{angles_py[2]:.2f}°", color="blue")

# Annotate side lengths at midpoints
midpoints = [
    ((A[0]+B[0])/2, (A[1]+B[1])/2, sides_py[0]),
    ((B[0]+C[0])/2, (B[1]+C[1])/2, sides_py[1]),
    ((C[0]+A[0])/2, (C[1]+A[1])/2, sides_py[2])
]
  \end{lstlisting}
\end{frame}

\begin{frame}[fragile]
\frametitle{Python Code}
\begin{lstlisting}
for (mx, my, d) in midpoints:
    plt.text(mx+0.2, my+0.2, f"{d:.2f}", fontsize=10, color="purple")

# Axis
plt.axhline(0, color="black", linewidth=0.5)
plt.axvline(0, color="black", linewidth=0.5)
plt.axis('equal')
plt.grid(True)
plt.title("Triangle ABC with Sides and Angles (from C)")
plt.legend()
plt.savefig("fig_3.png")
plt.show()
  \end{lstlisting}
\end{frame}

 \begin{frame}{Plot}
    \centering
    \includegraphics[width=\columnwidth, height=0.8\textheight, keepaspectratio]{../Figs/fig_3.png}     
\end{frame}

\end{document}





